%%--------------------------------------------------------------------------
%%  SECTION 5: FAST SPECTRAL ALGORITHMS
%%--------------------------------------------------------------------------
\section{Fast Spectral Algorithms for Kinetic Equations}

\noindent\textit{Why it matters.}\enspace
Classical fast algorithms play an indispensable dual role in my research
program.  They are the reference solvers that generate ground-truth data
for training and validating all machine learning models, and they are
independent scientific contributions that extend the computational frontier
of kinetic theory.  Without reliable, efficient classical solvers, the
benchmarking of any AI-based method is scientifically unsound.

The direct evaluation of the Boltzmann collision integral at velocity-space
resolution $N^3$ requires $O(N^6)$ operations per time step---completely
intractable for three-dimensional non-equilibrium gas dynamics even on
modern supercomputers.  Fast spectral methods that exploit the convolution
structure of the collision operator in Fourier space reduce this to
$O(N^4 \log N)$ for the elastic Boltzmann operator, but the \emph{inelastic}
Boltzmann operator---governing granular gases, vibrationally excited molecular
gases, and certain plasma-neutral interactions---had not been treated by any
algorithm achieving better than $O(N^6)$ prior to our work.

During my postdoctoral appointment at Purdue University (mentored by
Jingwei~Hu), I developed a \emph{fast spectral method for the inelastic
Boltzmann collision operator}~\cite{HU2019119}, achieving $O(N^2 \log N)$
complexity per velocity mode by uncovering a previously unrecognized
convolution structure in the inelastic collision kernel after a suitable
change of variables in Fourier space.  This is, to our knowledge, the
first sub-cubic algorithm for the inelastic operator.  The method handles
the full non-equilibrium collision dynamics and has been validated against
Direct Simulation Monte Carlo (DSMC) benchmarks for heated granular gas
flows across a broad range of inelasticity parameters.

The implementation is released as open-source software:
\begin{itemize}
  \item \href{https://github.com/mazhengcn/fsm-inelastic-boltzmann}{%
    \texttt{mazhengcn/fsm-inelastic-boltzmann}}: the complete spectral
    solver for the inelastic Boltzmann operator.
  \item \href{https://github.com/mazhengcn/kipack}{\texttt{mazhengcn/kipack}}:
    a high-performance Python package supporting both elastic and inelastic
    Boltzmann operators, with GPU acceleration via JAX/CuPy and a clean
    NumPy interface.
\end{itemize}
These tools are actively used in our group and by external collaborators to
generate training data for APNNs and DeepRTE, closing the loop between fast
classical algorithms and modern machine learning.

Earlier work~\cite{APUQ-Jin2017,APUQ-Jin2018} developed rigorous AP--UQ
methods for uncertain transport and hyperbolic equations via stochastic
Galerkin discretizations, establishing uniform spectral convergence from the
kinetic regime to the diffusion limit in the presence of random inputs.  These
results provided both the mathematical framework and the design intuition that
informed the APNN architecture, and they remain the most complete rigorous
analysis of uncertainty quantification for multiscale kinetic theory.
